\appendix
\chapter{Appendices}
\label{chap:appendix}

\section{Anomaly template CSV format}
\label{app:csv-format}

Each anomaly is encoded as a CSV file in
\texttt{data-gen/anomalies/}. The columns are described below.

\begin{table}[H]
    \centering
    \footnotesize
    \begin{tabular}{lp{0.62\linewidth}}
        \toprule
        Column & Meaning \\
        \midrule
        \texttt{time}    & Normalised time of the keypoint, in $[0, 1]$. \\
        \texttt{value}   & Normalised amplitude of the keypoint, in $[0, 1]$. \\
        \texttt{t\_minus}, \texttt{t\_plus} & Boolean flags allowing
        the keypoint to jitter in the negative or positive time direction. \\
        \texttt{v\_minus}, \texttt{v\_plus} & Same, for the amplitude axis. \\
        \texttt{interp}  & Interpolation mode for the segment that
        \emph{starts} at this keypoint:
        \texttt{linear} (straight line),
        \texttt{exp\_up} (concave-up exponential),
        \texttt{exp\_down} (concave-down exponential),
        \texttt{bell} (smoothed half-cosine). \\
        \bottomrule
    \end{tabular}
    \caption{Schema of the anomaly-template CSV files.}
    \label{tab:csv-schema}
\end{table}

\paragraph{Example: \texttt{bell.csv}.}
\begin{lstlisting}[caption={Encoding of the \texttt{bell} template.}, label={lst:bell-csv}]
time,value,t_minus,t_plus,v_minus,v_plus,interp
0.0,0.0,0,0,0,0,bell
0.5,1.0,1,1,1,1,bell
1.0,0.0,0,0,0,0,linear
\end{lstlisting}

The first and last keypoints have all jitter flags off, so the template
always starts and ends at $(0,0)$ and $(1,0)$. The middle keypoint can
jitter in both time and amplitude; the two half-cosine segments make
the curve smooth.

\section{Hyperparameter reference}
\label{app:hp}

\subsection{Data generator}
\label{app:hp-gen}

\begin{table}[H]
    \centering
    \footnotesize
    \begin{tabular}{lll}
        \toprule
        Group & Parameter & Value \\
        \midrule
        \multirow{2}{*}{Size}     & Total points          & $10^{7}$ \\
                                   & Chunk size            & $5 \cdot 10^{4}$ \\
        \midrule
        \multirow{2}{*}{Anomalies}& Number of injections  & uniform in $[5\,000, 8\,000]$ \\
                                   & Injection jitter      & $\pm 200$ samples \\
        \midrule
        \multirow{3}{*}{Shapes}   & Amplitude (in $\sigma$) & uniform in $[2.5, 7.6]$ \\
                                   & Period (samples)        & uniform in $[200, 800]$ \\
                                   & Deformation variance    & uniform in $[0.02, 0.10]$ \\
        \midrule
        \multirow{4}{*}{Background}& Global mean $\mu$    & $99.26$ \\
                                   & HF noise std $\sigma$ & $2.53$ \\
                                   & OU $\theta$           & $4 \cdot 10^{-6}$ \\
                                   & OU $\sigma_{\mathrm{OU}}$ & $0.0373$ \\
        \bottomrule
    \end{tabular}
    \caption{Configuration of \texttt{generate\_custom\_dataset.py} used
    to produce the dataset for this report.}
    \label{tab:gen-config}
\end{table}

\subsection{Training}
\label{app:hp-train}

\begin{table}[H]
    \centering
    \footnotesize
    \begin{tabular}{ll}
        \toprule
        Hyperparameter & Value \\
        \midrule
        Window size $W$                          & $512$ samples \\
        Sliding-window stride at training        & $32$ samples \\
        Sliding-window stride at inference       & $10$ samples \\
        Window labelling threshold               & $10\%$ anomaly points \\
        Per-window normalisation                 & subtract window mean, divide by global $\sigma_{\mathrm{train}}$ \\
        $\sigma_{\mathrm{train}}$ used in this report & $4.21$ \\
        Train/val/test split                     & temporal $80/20$ then random $80/20$ on train \\
        Optimizer                                & Adam, $\beta_1 = 0.9$, $\beta_2 = 0.999$, $\varepsilon = 10^{-8}$ \\
        Initial learning rate                    & $10^{-3}$ \\
        Batch size                               & $128$ \\
        Maximum epochs                           & $35$ \\
        LR scheduler                             & \texttt{ReduceLROnPlateau}(patience=$2$, factor=$0.5$) \\
        Early stopping                           & patience $10$, min delta $5 \cdot 10^{-3}$ \\
        Loss                                     & weighted categorical cross-entropy \\
        Class weights $\mathbf{w}$               & $(0.46, 1.18, 1.27, 1.32, 1.24, 1.15, 1.31)$ \\
        \bottomrule
    \end{tabular}
    \caption{Shared training configuration.}
    \label{tab:train-config}
\end{table}

\subsection{Inference and event extraction}
\label{app:hp-inf}

\begin{table}[H]
    \centering
    \footnotesize
    \begin{tabular}{ll}
        \toprule
        Hyperparameter & Value \\
        \midrule
        Sliding-window stride & $10$ samples \\
        Inference batch size  & $256$ windows \\
        Smoothing window      & $30$-sample rolling mean on the per-sample probability trace \\
        Background-probability threshold for an event & $\hat{p}_0 < 0.3$ \\
        Tolerance gap between adjacent events & $50$ samples \\
        \bottomrule
    \end{tabular}
    \caption{Configuration of the inference engine and event extractor.}
    \label{tab:inf-config}
\end{table}

\section{Reproducing the contents of this report}
\label{app:reproduce}

The repository layout is:
\begin{lstlisting}[basicstyle=\ttfamily\footnotesize, caption={Top-level project structure.}]
pole-anomaly/
  biblio/                  - the project brief and reference papers
  data-gen/                - real data + synthetic generator
    anomaly_generator.py     - template deformation + interpolation
    generate_custom_dataset.py - end-to-end dataset generator
    evaluate_real_data.py    - statistical characterisation of real data
    visualize_anomalies_templates.py
    interactive_labeling_helper.py
    anomalies/               - the six template CSV files
    data/                    - real recordings + generated CSV
    logs/                    - PNG/HTML inspection plots
  architectures/           - models, training, inference
    train.py                 - shared training pipeline (selects model)
    visualize_inference.py   - synthetic-test event-level evaluation
    inference_real_data.py   - real-data deployment trial
    robustness_test.py       - noise/deformation sweep
    inference_engine.py      - sliding-window + event matching
    model_selection.py       - model registry & CLI selector
    saved_models/            - 1d_cnn.py, resNet.py, cnn_attention.py
    logs/                    - per-model training/inference outputs
  documentation/           - this report
    main.tex, preamble.tex, sections/*.tex
    figures/build_figures.py - reproduces all PNGs in this report
\end{lstlisting}

The numbers and figures in this report can be reproduced from a clean
checkout with the following command sequence:

\begin{lstlisting}[basicstyle=\ttfamily\footnotesize]
# 1. Build the synthetic dataset (about 10 minutes; 600 MB on disk)
py data-gen/evaluate_real_data.py         # writes data/real_data_params.json
py data-gen/generate_custom_dataset.py    # writes data/synthetic_custom_dataset.csv

# 2. Make a copy of the dataset under architectures/data/ as
#    training_dataset.csv and validation_dataset.csv. (Same file twice is fine
#    if the user only wants to re-train; the temporal split inside train.py
#    keeps the test segment leakage-free.)

# 3. Train each architecture (~5-10 minutes each on an RTX-class GPU)
py architectures/train.py --model 1d_cnn
py architectures/train.py --model resnet
py architectures/train.py --model cnn_attention

# 4. Synthetic-test event-level evaluation
py architectures/visualize_inference.py --model 1d_cnn
py architectures/visualize_inference.py --model resnet
py architectures/visualize_inference.py --model cnn_attention

# 5. Robustness sweeps
py architectures/robustness_test.py --model 1d_cnn
py architectures/robustness_test.py --model resnet
py architectures/robustness_test.py --model cnn_attention

# 6. Real-data deployment trial
py architectures/inference_real_data.py --model 1d_cnn
py architectures/inference_real_data.py --model resnet
py architectures/inference_real_data.py --model cnn_attention

# 7. Rebuild the figures of this report
py documentation/figures/build_figures.py

# 8. Compile the report (with a local TeX install)
cd documentation && pdflatex main && pdflatex main
\end{lstlisting}

\section{Software environment}
\label{app:env}

All experiments were run on a single laptop with the following
software stack:
\begin{itemize}[leftmargin=*, itemsep=0.2em]
    \item Windows 11, Python $3.13$.
    \item PyTorch $2.12$ (CUDA 12.8), running on an NVIDIA RTX 5060
    Laptop GPU.
    \item NumPy, Pandas, Matplotlib, SciPy, scikit-learn for data
    handling, statistics and reporting.
    \item Plotly for the interactive HTML diagnostics shipped under
    \texttt{logs/}.
\end{itemize}
A pinned dependency list is included in the top-level
\texttt{requirements.txt} of the repository.
